%======================================================================
\section{Applying IHS to Multivalued Regular MaxSAT}
\label{sec:regular}
%======================================================================

%----------------------------------------------------------------------
\subsection{Where Regular MaxSAT Sits}
\label{sec:regular-place}
%----------------------------------------------------------------------

The preceding sections bracket a spectrum. \cref{sec:part1,sec:part2} treat
\emph{Boolean} MaxSAT, where each soft clause is either satisfied or
falsified, the falsification indicators $y_i$ are binary, and IHS solves a
minimum-cost hitting-set IP whose dual is the core-packing LP~\eqref{eq:pack}.
\cref{sec:part3} treats the other end, general Weighted CSP, where variables
range over finite domains, soft objects are multi-valued cost functions, and
the master becomes a minimum-cost \emph{hitting-vector} (MHV) problem over cost
vectors (\cref{sec:vecihs}). This section fills in the natural midpoint:
\emph{multivalued Regular MaxSAT}, in which variables are multi-valued but the
soft objects retain \emph{order} (regular) structure.

The question worth asking is not merely ``can IHS be run here''---it can, and
\cref{sec:regular-enc} says how in one paragraph---but \emph{what the master
becomes}. The answer organizes the whole section: the Benders/Dantzig--Wolfe
hitting-set machinery of \cref{sec:part1,sec:part2} survives multivaluedness
essentially intact. The regular master is the ordinary hitting-set IP of
\cref{alg:ihs} augmented by one \emph{ladder} of order inequalities per
variable; those ladders form a totally unimodular block, so they are
gap-neutral (\cref{prop:gapneutral}). The general MHV master of \cref{sec:part3}
is what remains once the cost functions \emph{lose} the order structure and the
ladder collapse no longer applies. Regular MaxSAT is thus the setting in which
multivaluedness is added to Boolean MaxSAT \emph{without} leaving the
duality square of \cref{sec:unified}.

%----------------------------------------------------------------------
\subsection{Multivalued Regular MaxSAT}
\label{sec:regular-def}
%----------------------------------------------------------------------

Fix variables $x_1,\dots,x_n$, each ranging over a finite totally ordered
domain; without loss of generality $x_j \in D_j = \{0,1,\dots,d_j-1\}$. A
\emph{regular literal} is a threshold constraint on a single variable, either
an up-literal $(x_j \ge a)$ or a down-literal $(x_j \le a)$ for some
$a \in D_j$; an assignment $I : x_j \mapsto D_j$ satisfies $(x_j \ge a)$ iff
$I(x_j) \ge a$, and dually for $(x_j \le a)$. A \emph{regular clause} is a
disjunction of regular literals, and a \emph{regular CNF formula} is a
conjunction of regular clauses. These are the ``regular'' signed CNF formulas
of many-valued logic, the fragment of signed CNF in which every sign is an
up-set or a down-set of the truth-value order~\citep{beckert2000signed}.

A \emph{weighted partial regular MaxSAT} instance $\langle \mathcal{H},
\mathcal{S}, w\rangle$ has hard regular clauses~$\mathcal{H}$, soft regular
clauses $\mathcal{S} = \{s_1,\dots,s_m\}$, and weights
$w : \mathcal{S} \to \mathbb{Q}_{>0}$; the goal is an assignment satisfying all
of $\mathcal{H}$ that minimizes the total weight of falsified soft clauses.
Setting every $d_j = 2$ recovers Boolean MaxSAT exactly. Regular MaxSAT is the
signed-logic face of WCSP: the translation of~\citet{ansotegui2007logic} maps a
WCSP with order-structured (in particular convex) cost functions onto exactly
this fragment, so the setting is neither exotic nor a toy---it is the standard
bridge between the two frameworks the memo studies.

%----------------------------------------------------------------------
\subsection{The Order Encoding and the Shallow Transfer}
\label{sec:regular-enc}
%----------------------------------------------------------------------

The quickest way to run IHS on a regular instance is to Booleanize it with the
\emph{order} (a.k.a.\ regular, ladder) encoding of finite-domain
variables~\citep{ansotegui2004mapping}. For each variable $x_j$ introduce
Boolean ladder variables $X_j^a$ for $a = 1,\dots,d_j-1$ with intended meaning
$X_j^a \equiv (x_j \ge a)$, fix $X_j^0 = \top$ and $X_j^{d_j} = \bot$, and add
the \emph{monotonicity axioms}
\begin{equation}
X_j^{a+1} \to X_j^{a}, \qquad\text{i.e.}\qquad (\neg X_j^{a+1} \vee X_j^{a}),
\qquad a = 1,\dots,d_j-2,
\label{eq:ladder}
\end{equation}
as \emph{hard} clauses. Under~\eqref{eq:ladder} every Boolean model corresponds
to a unique domain value $I(x_j) = \max\{a : X_j^a = 1\}$, and every regular
literal becomes an ordinary Boolean literal: $(x_j \ge a) \mapsto X_j^a$ and
$(x_j \le a) \mapsto \neg X_j^{a+1}$. A regular clause therefore Booleanizes to
an ordinary clause over the ladder variables.

The consequence is that \emph{the entire apparatus of \cref{sec:part1,sec:part2}
transfers verbatim}. Cores are subsets of soft (regular) clauses unsatisfiable
together with $\mathcal{H}$ and the monotonicity axioms; the master is the
hitting-set IP~\eqref{eq:hit}; the SAT oracle is a call on the order-encoded
formula, or a native regular / many-valued SAT solver~\citep{bejar2007regular},
or any finite-domain CP/SMT backend. The Benders reading (core-guided,
\cref{sec:benders}) and the Dantzig--Wolfe reading (IHS, \cref{sec:dw}) apply
unchanged, because the only thing that changed is \emph{how the subproblem
decides satisfiability}, not the structure of the master or the cut. This is the
honest one-line answer to ``how do I apply IHS to regular MaxSAT.'' The rest of
the section is about what is \emph{gained} by not throwing the order structure
away at the encoding step.

%----------------------------------------------------------------------
\subsection{Multivalued Costs as Soft Regular Literals}
\label{sec:regular-costs}
%----------------------------------------------------------------------

The order structure earns its keep as soon as costs are multi-valued. A unary
cost function $f_j : D_j \to \mathbb{Q}_{\ge 0}$ that is non-decreasing (the
common WCSP case: deviation penalties, tardiness, resource overuse) is
represented by soft regular \emph{down}-literals with telescoped weights. For
each threshold $a = 1,\dots,d_j-1$ introduce
\begin{equation}
\ell_j^{a} \;:=\; (x_j \le a-1) \;\equiv\; (x_j < a),
\qquad w(\ell_j^{a}) \;=\; \Delta_j^{a} \;:=\; f_j(a) - f_j(a-1) \;\ge\; 0.
\label{eq:staircase}
\end{equation}
Write $y_j^{a} := \mathbf 1[\ell_j^{a}\text{ falsified}] = \mathbf 1[x_j \ge a]$.
If the assignment sets $x_j = v$, the falsified literals are exactly
$\{\ell_j^{a} : a \le v\}$, of total weight
$\sum_{a=1}^{v}\Delta_j^{a} = f_j(v) - f_j(0)$. So, up to the additive constant
$\sum_j f_j(0)$, the falsification cost equals $\sum_j f_j(x_j)$: the staircase
telescopes. (A non-monotone $f_j$ is handled by splitting it into monotone
pieces, or with up-literals as well; the clean single-direction staircase is
the convex case.)

Two facts make this more than a re-encoding. First, because the falsified
literals of a variable are a \emph{prefix} of its threshold chain, the
indicators are monotone: $y_j^{a} \ge y_j^{a+1}$. These are exactly the
primal-side image of the monotonicity axioms~\eqref{eq:ladder}, now expressed on
the \emph{master} variables. Second, a core acquires a \emph{threshold}
structure. A core is a set of down-literal assumptions
$\{x_j \le a_j-1 : j \in J\}$ that is unsatisfiable with $\mathcal{H}$; it
certifies $\mathcal{H} \models \bigvee_{j \in J}(x_j \ge a_j)$, a genuine
regular clause, whose linearization is the single-indicator-per-variable
hitting-set row $\sum_{j\in J} y_j^{a_j} \ge 1$. This is precisely the cost-vector
core of \cref{sec:vecihs} specialized to unary regular costs: the induced CSP
$F_{\vec v}$ there is here a conjunction of hard upper-bound literals
$\{x_j \le a_j-1\}$, and its unsatisfiability is a regular clause.

Collecting the objective~\eqref{eq:staircase}, the covering rows, and the ladder
rows gives the \textbf{regular hitting-set master}
\begin{align}
\min \;\; & \sum_{j} \sum_{a=1}^{d_j-1} \Delta_j^{a}\, y_j^{a}
   \label{eq:reg-master-obj}\\
\text{s.t.}\;\; & \sum_{j \in J_t} y_j^{a_j^{t}} \ge 1
   && \forall\text{ cores } t, \label{eq:reg-master-cover}\\
& y_j^{a} \ge y_j^{a+1}
   && \forall j,\ a = 1,\dots,d_j-2, \label{eq:reg-master-ladder}\\
& y_j^{a} \in \{0,1\}. \notag
\end{align}
Solving~\eqref{eq:reg-master-ladder} together with the objective, choosing the
top indicator $y_j^{a_j}=1$ forces $y_j^{a}=1$ for all $a \le a_j$ and hence
pays $\sum_{a\le a_j}\Delta_j^{a} = f_j(a_j)-f_j(0)$: the true cost of lifting
$x_j$ to $a_j$. The master is still a $0/1$ program over soft-clause
indicators---\emph{not} the general MHV program of \cref{sec:vecihs}---with the
multivaluedness absorbed into (i) which threshold indicator a core selects and
(ii) the ladder rows that make the objective charge the whole prefix.

\paragraph{A small worked instance.} Take $x_1, x_2 \in \{0,1,2\}$ with
$f_1(v)=f_2(v)=v$, so each variable contributes soft literals
$\ell_j^{1}=(x_j<1)$ and $\ell_j^{2}=(x_j<2)$ of weight $1$, plus the ladder row
$y_j^{1} \ge y_j^{2}$. Let the single hard clause be the regular clause
$\mathcal{H} = \{(x_1 \ge 2)\vee(x_2 \ge 2)\}$ (``some variable must reach its
top''). The oracle on $\mathcal{H}\cup\{x_1\le1,\ x_2\le1\}$ is UNSAT and returns
the core $\{\ell_1^{2},\ell_2^{2}\}$, cut $y_1^{2}+y_2^{2}\ge1$. The
master~\eqref{eq:reg-master-obj}--\eqref{eq:reg-master-ladder} then sets, e.g.,
$y_1^{2}=1$, forcing $y_1^{1}=1$ by the ladder, for cost $2 = z^\star$ (lift
$x_1$ to $2$); the next oracle call on $\mathcal{H}\cup\{x_2\le1\}$ succeeds
($x_1=2, x_2=0$). Note that here $z^\star = z^\star_{\mathrm{LP}} = 2$: the
multivaluedness alone creates no gap. To reproduce a $K_3$-style gap
(\cref{sec:counter}) one needs an odd cycle \emph{among variables}, not within a
single variable's ladder---which is exactly the content of the next subsection.

%----------------------------------------------------------------------
\subsection{What Multivaluedness Costs the Master: Gap-Neutrality}
\label{sec:regular-gap}
%----------------------------------------------------------------------

The pressing question for the equivalence analysis of \cref{sec:part2} is
whether multivaluedness enlarges the integrality gap of the hitting-set
polytope---the exact quantity that separates pure Dantzig--Wolfe (IHS-LP) from
practical IHS-IP (\cref{prop:nonequiv}). It does not.

\begin{proposition}[Ladder rows are gap-neutral]
\label[proposition]{prop:gapneutral}
In the regular master~\eqref{eq:reg-master-obj}--\eqref{eq:reg-master-ladder},
the ladder constraints~\eqref{eq:reg-master-ladder} form a totally unimodular
block, and they are valid for every integer-feasible point. Consequently,
adding them to the LP relaxation of the master weakly \emph{lowers} the
integrality gap and never raises it: the LP value weakly increases toward
$z^\star$ while the IP value is unchanged. In particular, any positive gap of
the regular master is inherited entirely from the covering
rows~\eqref{eq:reg-master-cover}, i.e.\ from conflict structure \emph{among
variables}, exactly as in the Boolean case.
\end{proposition}

\begin{proof}
Each row of~\eqref{eq:reg-master-ladder} is $y_j^{a} - y_j^{a+1} \ge 0$, with a
single $+1$ and a single $-1$; per variable these form an interval (difference)
matrix, and across variables the block is block-diagonal, hence totally
unimodular. Validity is immediate: any assignment sets a prefix of each ladder
to $1$, so $y_j^{a}\ge y_j^{a+1}$ holds at every integer point, and the IP
optimum is unchanged. Since the rows are valid, they cut off only fractional
points; adjoining them shrinks the LP-feasible region, so the minimization LP
value weakly increases. With $z^\star$ fixed, $z^\star - z^\star_{\mathrm{LP}}$
weakly decreases. The residual gap is then a property of the covering block
alone, which is the Boolean hitting-set matrix of~\eqref{eq:hit}.
\end{proof}

Two readings follow. Negatively, multivaluedness is ``free'' for the
duality story: the exact-equivalence proposition \cref{prop:equiv} (IHS-LP $=$
DW) and the non-equivalence proposition \cref{prop:nonequiv} (integer master
strictly stronger) both carry over to regular MaxSAT unchanged, with the ladder
rows as gap-neutral valid inequalities and the $K_3$ triangle of
\cref{sec:counter} still the minimal witness of a gap. Positively, the ladder
rows can only \emph{help}: on instances where the order structure interacts with
the conflict structure they may close part of the gap for free---the same
phenomenon by which abstract cores strengthen the LP master in
\cref{sec:abstract-cores}, obtained here from the encoding rather than from an
added cardinality constraint.

\begin{remark}[Regular literals are ordered counting literals]
\label[remark]{rem:regular-abstract}
There is a precise kinship between the threshold indicators $y_j^{a}$ and the
counting literals of abstract cores (\cref{sec:abstract-cores}). An abstract
core introduces a literal asserting ``at least $k$ clauses of $A$ are
falsified,'' inducing the rank-$k$ master row $\sum_{s_i \in A} y_i \ge k$ of
\eqref{eq:abstract-row} over an \emph{unordered, interchangeable} set $A$. A
regular variable ships with exactly such counting literals already---$(x_j \ge
a)$ is ``at least $a$ of the ladder is falsified''---but its clauses are
\emph{totally ordered}, and the ladder rows~\eqref{eq:reg-master-ladder} are what
encode that order (a nested prefix) rather than a symmetric count. Abstract
cores, viewed through this lens, \emph{re-introduce} on flat Boolean softs the
threshold structure that regular MaxSAT carries natively; conversely,
aggregating a regular variable's ladder into a single abstract core is
degenerate, since its thresholds already are the counts. The aggregation lever
of \cref{sec:abstract-cores} and the encoding of \cref{sec:regular-costs} are
two routes to the same rank-$k$ rows.
\end{remark}

%----------------------------------------------------------------------
\subsection{The Oracle and Core Improvement}
\label{sec:regular-oracle}
%----------------------------------------------------------------------

Pricing---the SAT oracle in the Dantzig--Wolfe reading of \cref{sec:dw}---has a
convenient form here. The candidate correction supplied by the master is a
vector of thresholds, one per variable: $H^\star$ is read off as ``leave $x_j$
free up to the largest $a$ with $y_j^{a}=0$.'' The oracle then decides
$\mathcal{H} \cup \{x_j \le a_j-1\}_j$, which an assumption-based solver
implements by asserting the ladder literals $\neg X_j^{a_j}$; an UNSAT answer
returns a core as a threshold vector, and the incrementality that makes ``lazy
cores'' free in \cref{sec:part3} works identically. Multi-core and disjoint-core
extraction (the transferred technique that dominated the ablations of
\cref{sec:part3}, \cref{fig:benchmarks}) carry over without change.

Core improvement (\textsc{ImprCore}, \cref{alg:ihs-lb}) also specializes
cleanly. A \emph{maximal} regular core is obtained by pushing each threshold
$a_j$ as high as the core condition permits---raising the enforced upper bound
$x_j \le a_j-1$ until satisfiability is regained---which is exactly the
component-wise increment of \cref{sec:alts} restricted to the ordered domain.
The ``prefer the lowest current value'' rule of \cref{sec:part3} becomes
``prefer to raise the variable whose threshold sits lowest on its ladder,''
making cores costlier to hit and so accelerating the lower bound. In short, the
Axis-2 design space of \citet{petrova2025empirical} instantiates on regular
instances with thresholds playing the role of cost-vector components, and the
Axis-1 hitting-vector relaxations (cost-bounded, greedy) apply to the regular
master~\eqref{eq:reg-master-obj}--\eqref{eq:reg-master-ladder} as written.

%----------------------------------------------------------------------
\subsection{The Spectrum: Boolean \texorpdfstring{$\subseteq$}{<=} Regular \texorpdfstring{$\subseteq$}{<=} WCSP}
\label{sec:regular-spectrum}
%----------------------------------------------------------------------

\cref{tab:spectrum} places the three settings on one axis. Reading down the
``Master'' column tells the whole story: the hitting-set IP of Boolean MaxSAT
gains, in the regular case, one totally unimodular ladder per variable (a
gap-neutral addition, \cref{prop:gapneutral}), and gains, in the WCSP case, the
full multi-valued hitting-vector program once the cost functions are no longer
order-structured and the ladder collapse of \cref{sec:regular-costs} no longer
applies.

\begin{table}[htbp]
\centering
\small
\setlength{\tabcolsep}{5pt}
\caption{Multivalued Regular MaxSAT as the structured midpoint between the
Boolean MaxSAT of \cref{sec:part1,sec:part2} and the general WCSP of
\cref{sec:part3}. The Dantzig--Wolfe / Benders duality of the memo survives
into the regular row intact; the WCSP row is what remains once the order
structure is dropped.}
\label{tab:spectrum}
\begin{tabular}{@{}llll@{}}
\toprule
\textbf{Setting} & \textbf{Variables} & \textbf{Soft object} & \textbf{Master} \\
\midrule
Boolean MaxSAT & binary & clause, indicator $y_i$ & hitting-set IP~\eqref{eq:hit} \\
Regular MaxSAT & multivalued, ordered & threshold literal $y_j^{a}$ & hitting-set IP $+$ ladders~\eqref{eq:reg-master-ladder} \\
WCSP & multivalued & cost function, value $v_i$ & min-cost hitting vector (MHV) \\
\bottomrule
\end{tabular}
\end{table}

This diagram also explains, in the memo's own terms, the empirical suitability
proxy that \cref{sec:part3} distilled from \citet{petrova2025empirical}: IHS
does best on WCSP instances with \emph{small domains and few distinct costs}.
Those are precisely the instances whose order encoding stays compact---few
ladder rungs per variable, hence a regular master~\eqref{eq:reg-master-obj}--%
\eqref{eq:reg-master-ladder} that departs only mildly from the Boolean
hitting-set IP on which the Dantzig--Wolfe analysis is tightest. Large domains
with many distinct costs (CELAR, in \cref{tab:wcspclasses}) push the instance
away from the regular midpoint and toward the full MHV master, where the clean
row--column duality of \cref{sec:dw} is diluted. The regular view thus does more
than add a setting: it locates \emph{why} the hitting-set approach tracks the
domain and cost granularity, which the Boolean and WCSP endpoints, taken alone,
leave unexplained.
